\hl{The work relevant to this thesis comes from research streams that
have developed largely separately: Bangla dialect processing, Bangla
question answering, Bangla language models, and applications of language
technology in healthcare and education. This chapter first introduces
the concepts the thesis builds on, then reviews the related work,
highlighting its achievements and limitations, and ends with the gaps
that the proposed framework addresses.}

\section{Preliminaries}\label{preliminaries}

\subsection{Bangla and Its Regional Varieties}\label{bangla-and-its-regional-varieties}

\hl{Standard Bangla (cholito bhasha) is the form used in formal writing,
education, broadcasting and administration. The regional varieties form
a continuum across Bangladesh and differ from the standard in
pronunciation, vocabulary, verb morphology and sometimes word order, in
some regions vastly so \cite{jawad2026fivedialects}, \cite{faria2023vashantor}. Chittagonian, spoken in the
south-east, and Sylheti, spoken in the north-east, are among the most
divergent. Online, dialect speakers write in Bangla script with phonetic
spellings of dialect words, or in Latin script (romanized Bangla, often
called Banglish); Sylheti also has a historical script of its own,
Sylheti Nagri \cite{jawad2026fivedialects}. A dialect sentence and its Standard equivalent
can share very few word forms, sometimes none, although they convey
exactly the same meaning, for example Sylheti \emph{kita} for Standard
\emph{ki} (``what''), which is why systems trained on Standard Bangla
misread dialect input.}

\subsection{Dialect Identification and Normalization}\label{dialect-identification-and-normalization}

\hl{Dialect identification predicts the variety a text is written in and
is usually treated as text classification \cite{faria2023vashantor}. Dialect normalization
rewrites a dialect text in the standard variety while keeping its
meaning, which makes it a form of machine translation \cite{faria2023vashantor}, \cite{jawad2026fivedialects}.
Its quality is measured with BLEU \cite{papineni2002bleu}, the character n-gram
F-scores chrF and chrF++ \cite{popovic2015chrf}, \cite{popovic2017chrfpp}, and word or character error
rate (WER, CER). A useful sanity check is the copy-input baseline, which
scores the unchanged dialect input against the Standard reference: a
system that scores below it has made the text worse \cite{jawad2026fivedialects}.}

Both tasks are long-standing problems in NLP for closely related
languages, varieties and dialects \cite{zampieri2020survey}. Dialect identification
is usually framed as supervised text classification, and in shared tasks
for Arabic it has reached high scores: the best systems of NADI 2023
reached a macro-F1 of 87.27 on country-level Arabic dialect
identification \cite{abdulmageed2023nadi}. Normalization is much harder than
identification. In the same shared task, the best systems that
translated Arabic dialects into Modern Standard Arabic reached only
14.76 and 21.10 BLEU on the two machine-translation subtasks
\cite{abdulmageed2023nadi}. Normalization is easier to learn when parallel
dialect--standard data exist, as in the MADAR corpus of 25 city
dialects \cite{bouamor2018madar}, and character-level models suit it because dialect
words are often spelling variants of standard words: for Swiss German,
character-based neural translation was the best normalization strategy
\cite{honnet2018swissgerman}.

Romanized input is a related problem. Bangla written in Latin letters
has no fixed spelling, so the same word appears in many forms, and
converting it back into Bangla script (back-transliteration) is a task
of its own. BanglaTLit provides 42.7k romanized--Bangla-script pairs for
this task, together with a 245.7k-sample romanized pre-training corpus
\cite{fahim2024banglatlit}. In this thesis, romanized questions are
normalized by the same expert that normalizes dialect questions, so
back-transliteration and dialect normalization are handled in one step.

\subsection{Question Answering}\label{question-answering}

\hl{In extractive QA the answer is a span of a given passage, scored
with exact match (EM) and token-level F1 \cite{ekram2022banglarqa}; in multiple-choice QA
the system picks an option and is scored by accuracy \cite{shafayat2024benqa}, \cite{shifa2025somajgyaan};
in open-domain QA the system must first find the evidence itself,
usually with a retriever followed by a reader \cite{karpukhin2020dpr}. When the answer
is generated in the system's own words, evaluation relies on semantic
similarity measures such as BERTScore \cite{zhang2020bertscore} and on checking whether
the generated answer is supported by the evidence \cite{es2024ragas}. A realistic
system must also recognise questions it cannot answer. SQuAD 2.0 added
more than 50,000 unanswerable questions written adversarially by crowd
workers to look answerable, and showed that models trained only on
answerable questions answer them confidently \cite{rajpurkar2018squad2}. Selective QA
trains a calibrator that lets a model abstain when it is unsure, trading
coverage for accuracy \cite{kamath2020selective}.}

\hl{Most of the early Bangla QA systems follow the reading-comprehension
setting, where the passage is given for context and answer searching. A
practical system must instead first retrieve relevant documents from a
knowledge base with a retriever and then pass the top-k documents to a
reader model, which is the retriever--reader design of open-domain QA
\cite{karpukhin2020dpr}. The framework in this research follows that design, with an
LLM as the reader.}

\subsection{Large Language Models and Parameter-Efficient Fine-Tuning}\label{large-language-models-and-parameter-efficient-fine-tuning}

\hl{Modern LLMs are Transformers pre-trained on large text collections
and then instruction-tuned. Updating all their parameters is too
expensive for most academic budgets. LoRA freezes each pretrained weight
matrix \emph{W} \ensuremath{\in} \ensuremath{\mathbb{R}}\textsuperscript{d\ensuremath{\times}k} and learns a low-rank update
\ensuremath{\Delta}\emph{W} = \emph{BA} with \emph{B} \ensuremath{\in} \ensuremath{\mathbb{R}}\textsuperscript{d\ensuremath{\times}r}, \emph{A} \ensuremath{\in}
\ensuremath{\mathbb{R}}\textsuperscript{r\ensuremath{\times}k} and \emph{r} \ensuremath{\ll} min(\emph{d}, \emph{k}) \cite{hu2022lora}.
Only \emph{A} and \emph{B} are trained, which reduces the number of
trainable parameters by several orders of magnitude, and different
adapters can be attached to the same base model for different tasks.
QLoRA stores the frozen base model in 4-bit NormalFloat (NF4) precision
with double quantization and paged optimizers and back-propagates
through it into the LoRA adapters, so billion-parameter models can be
fine-tuned on one GPU with little loss in quality \cite{dettmers2023qlora}.}

\subsection{Retrieval-Augmented Generation}\label{retrieval-augmented-generation}

\hl{RAG combines a retriever that finds passages relevant to the
question with a generator that writes the answer from them \cite{lewis2020rag}.
Sparse retrieval such as BM25 scores documents by weighted term overlap
and works well when question and document share words \cite{robertson2009bm25}. Dense
retrieval encodes questions and passages as vectors and ranks by vector
similarity, which also matches paraphrases and synonyms \cite{karpukhin2020dpr};
multilingual embedding models such as multilingual E5 and BGE-M3 support
Bangla \cite{chen2024bgem3}, \cite{wang2024me5}. Hybrid systems fuse sparse and dense
rankings, typically with reciprocal rank fusion (RRF) \cite{cormack2009rrf}. Self-RAG
trains the model to decide when to retrieve and to critique both the
retrieved passages and its own output \cite{asai2024selfrag}, and Adaptive-RAG uses a
small classifier to choose between no retrieval, single-step and
multi-step retrieval according to question complexity \cite{jeong2024adaptiverag}.}

\subsection{Mixture of Experts and Routing}\label{mixture-of-experts-and-routing}

\hl{A mixture of experts is a set of specialised models with a gating
function that decides how much each model contributes to each input
\cite{jacobs1991adaptive}. In current LLMs, sparse mixture-of-experts layers send each
token to a few expert sub-networks inside one model, which increases
capacity at little extra compute \cite{shazeer2017moe}, \cite{fedus2022switch}. This thesis applies
the idea at the system level rather than inside the network: separate
lightweight routers send each question to one normalization expert and
one domain expert, and each expert is an independently trained LoRA
adapter on a shared base model. This is closer to Adaptive-RAG's small
routing classifier \cite{jeong2024adaptiverag} and to adapter-based transfer \cite{houlsby2019adapters} than
to a sparse MoE layer, so we call it a \emph{Mixture-of-Adapters}. It
keeps each expert small, interpretable and replaceable.} Related methods
that combine several adapters are reviewed in
Section~\ref{combining-and-routing-lora-adapters}.

\section{Review of Existing Research}\label{review-of-existing-research}

\hl{Research relevant to dialect-aware Bangla QA has expanded rapidly
since 2022. The reviewed work falls into nine themes: dialect resources
and translation; dialects in health and LLM prompting; Bangla QA
benchmarks; Bangla-capable language models; Bangla healthcare and
education applications; dialect processing beyond Bangla; combining and
routing LoRA adapters; re-ranking and hybrid retrieval; and grounding,
abstention and dialect robustness. Table~\ref{tab:related-studies} summarises the most closely related studies.}

\subsection{Dialect Resources and Dialect--Standard Translation}\label{dialect-resources-and-dialectstandard-translation}

\hl{Vashantor was the first large benchmark for translating Bangla
regional dialects into Standard Bangla \cite{faria2023vashantor}. It is multilingual
(Bangla, romanized Bangla and English) and covers five regions
(Chittagong, Noakhali, Sylhet, Barishal and Mymensingh) with 32,500
sentences: 2,500 Standard sentences rendered in every region give 12,500
dialect--Standard pairs, with an official split of 1,875 training, 250
validation and 375 test sentences per region. In the first version of
the paper, fine-tuned BanglaT5 reached a dialect\ensuremath{\rightarrow}Standard BLEU of 44.03
for Chittagong and 51.08 for Sylhet (mT5: 36.75 and 51.32), and region
detection reached 85.86\% accuracy, with Sylhet the hardest region (F1
0.75) \cite{faria2023vashantor}. The revised version introduces DialectBanglaT5 and
DialectBanglaBERT, reporting a best BLEU of 71.93 and lowest WER of
0.1470 (Mymensingh) and 89.02\% region-detection accuracy (Chittagong F1
0.9241). These results show that fine-tuning works across dialects, but
the data are short everyday sentences, not domain-specific, and
translation is the end task.}

\hl{Several newer corpora extend the dialect data. Kothon provides
human-validated Chittagonian--Standard and Sylheti--Standard sentence
pairs \cite{kothon2026}, ONUBAD covers several regional dialects \cite{sultana2025onubad}, and
ChatgaiyyaAlap provides Chittagonian--Standard pairs \cite{chatgaiyyaalap2025}. These
corpora supply most of the training pairs of our normalization expert.
5-Dialects-BN contains 6,000 manually verified utterances in five
dialects (1,900 Chittagong, 1,500 Noakhali, 1,200 Sylhet, 700 Barisal,
700 Rangpur), each with a romanized transliteration, a Standard Bangla
translation, an English translation and a subjectivity label; Cohen's \ensuremath{\kappa}
for the Standard Bangla renderings was 0.74 \cite{jawad2026fivedialects}. It evaluated seven
LLMs with zero-shot, few-shot and chain-of-thought prompting and with
LoRA. LoRA (\emph{r} = 16, \emph{\ensuremath{\alpha}} = 32) on 160 examples per dialect
raised Mistral-7B from 16.2 to 73.6 BLEU on dialect\ensuremath{\rightarrow}English translation,
above the zero-shot score of Gemini 3 Flash (46.4). Dialect\ensuremath{\rightarrow}Standard
normalization was evaluated only with prompting: some open models scored
below the copy-input baseline (39.66 chrF++), a fine-tuned BanglaT5
reached 44.47 chrF++, and the best system, few-shot Gemini, reached
78.70 chrF++. All models degraded on romanized input.}

\hl{Work in the opposite direction translates Standard Bangla into
dialects. Bridging Dialects fine-tuned BanglaT5, mT5 and mBART-50 to
translate Standard Bangla into the five Vashantor dialects; after 50
epochs BanglaT5 reached a WER of 0.2380 and CER of 0.1216 on Mymensingh
but 0.7066 and 0.3706 on Chittagong, with Chittagong and Sylhet the
hardest \cite{khandaker2024bridging}. BanglaDialecto built an end-to-end speech
standardization pipeline for Noakhali: about 10 hours of speech, a
fine-tuned Whisper recogniser (CER 0.8\%, WER 1.5\%), BanglaT5
translation into Standard Bangla (BLEU 41.6) and speech synthesis
\cite{samin2024bangladialecto}. It shows that dialect standardization can serve as an interface
for other systems, but it covers one dialect and stops before any
downstream task. None of these works connects dialect translation to
retrieval or to answering a user's question.}

\subsection{Dialects in Health and LLM Prompting}\label{dialects-in-health-and-llm-prompting}

\hl{BanglaCHQ-Prantik translated the 2,350 consumer health questions of
BanglaCHQ-Summ into Sylheti and 500 of them into Chittagonian with
native-speaker translators, creating the first medical-domain dialect
benchmark \cite{prantik2025}. The smaller Chittagonian subset reflects the fact
that few annotators can read and write Chittagonian and that it has no
standard orthography. It evaluated Qwen 2.5 3B, Gemma 3 1B, GPT-4o-mini
and Gemini 2.5 Flash with zero-shot, one-shot, five-shot and
chain-of-thought prompting in the Standard\ensuremath{\rightarrow}dialect direction; the best
BLEU was 23.67 (Sylheti, Gemini, five-shot) against 14.78 for
Chittagonian, with frequent errors in medical terms, numerals, idioms
and spelling. This benchmark contains exactly the kind of dialectal
health questions a healthcare QA system must handle, and it shows that
prompting frontier models is not enough. Sylheti-CAP studied Sylheti
translation in both directions on the Vashantor test set with
context-aware prompting of several LLMs, reaching at most 53.72 chrF for
Sylheti\ensuremath{\rightarrow}Standard \cite{prama2025sylheticap}. A comparison of RAG pipelines for
Standard\ensuremath{\rightarrow}dialect translation in six dialects found that ten LLMs reached
only 5--12 BLEU (WER 67--84\%) without retrieved examples, and that
retrieving aligned sentence pairs worked better than retrieving
transcripts (Chittagong BLEU 9\ensuremath{\rightarrow}26, WER 76\%\ensuremath{\rightarrow}55\%) \cite{sami2025comparative}. A smaller
model, Llama-3.1-8B (WER 51.18), even outperformed the much larger
GPT-OSS-120B (WER 52.65) \cite{sami2025comparative}, which indicates that parallel
examples matter more than model size. The same group benchmarked 19 LLMs
on questions in nine Bengali dialects and found far lower quality for
Chittagonian (5.44/10) than for Tangail (7.68/10) \cite{sami2026dialectbias}. AmarDoctor,
a multilingual triage application (81.08\% top-1 diagnostic precision on
185 clinical vignettes), took a lexicon approach: its medical team
linked 908 symptoms to 3,237 local variants, including Chittagonian and
Sylheti, and the authors report that a general speech-to-text service
struggled with dialect speech because of phonetic differences \cite{nahar2025amardoctor}. A
lexicon, however, only covers the terms it contains. Together, these
works show that prompting alone is weak for dialects, that retrieved
examples help, and that medical dialect data are scarce, especially for
Chittagonian.}

\subsection{Bangla Question Answering Benchmarks}\label{bangla-question-answering-benchmarks}

\hl{BanglaRQA is a human-annotated reading-comprehension dataset of
14,889 question--answer pairs on 3,000 Bangla Wikipedia passages, 3,631
of them unanswerable, with four question types (factoid, causal,
confirmation and list) and three answer types (single span, multiple
spans and yes/no); fine-tuned BanglaT5 reached 62.42 EM and 78.11 F1,
with weaker performance on the more difficult question types \cite{ekram2022banglarqa}.
The passages cover a wide variety of general-knowledge topics.}

\hl{NCTB-QA contains 87,805 question--answer pairs from 50
national-curriculum textbooks (grades 1--10), 57.25\% answerable and
42.75\% unanswerable with plausible distractors, split 80/10/10 so that
no context appears in two splits; fine-tuning raised BERT from 0.150 to
0.620 F1 \cite{nctbqa}. The dataset is closely aligned with the national
curriculum, which makes it a natural basis for the education expert;
however, its questions were model-generated from the textbook contexts,
so real student question types may be under-represented.}

\hl{BEnQA has 5,161 parallel Bangla--English science multiple-choice
questions from grade 8, 10 and 12 board exams; LLMs were consistently
less accurate in Bangla, chain-of-thought helped reasoning questions,
and appending the English translation of the question helped in all
subjects \cite{shafayat2024benqa}. SOMAJGYAAN has 4,234 multiple-choice questions on
Bangladeshi history, geography, economics, law and society, extracted
from national-curriculum books, where closed models led (Gemini 2.0
Flash 0.73, GPT-4o 0.69) and LoRA fine-tuning of open models did not
close the gap \cite{shifa2025somajgyaan}.}

\hl{Earlier systems were extractive and assumed that the correct passage
is given. Mayeesha et al. translated SQuAD 2.0 into Bangla, fine-tuned
multilingual BERT, DistilBERT and RoBERTa, tested them on a small
human-annotated set from Bangla Wikipedia, and compared the model
answers with answers from Bangla-speaking children as a human reference
point \cite{mayeesha2021bengaliqa}. Khondoker et al. built about 3,000 question--answer
pairs from NCTB textbooks of grades 6--10, where Bangla-BERT reached F1
0.75 and EM 0.53 and RoBERTa Base performed worst (F1 0.19); dataset
size and spelling errors were the main constraints \cite{khondoker2024unlocking}.}

\hl{BanglaCQA extends BanglaRQA with 6.3K counterfactual contexts
(42.28\% more contexts than the original) to test whether models use the
passage or their memory; under counterfactual passages, models largely
followed the passage and could not recover memorised answers (parametric
similarity about 0.13--0.14) \cite{azmary2026banglacqa}. Combining factual, counterfactual
and answerability data improved contextual answering, chain-of-thought
prompting helped decoder-only LLMs distinguish factual from
counterfactual knowledge, and the authors found LLM-based semantic
similarity more reliable than traditional overlap scores for comparing
Bangla answers \cite{azmary2026banglacqa}. For a RAG system, this means that the generator
has to be trained to rely on the retrieved evidence. All of these
benchmarks are in Standard Bangla, and none contains dialect questions.}

\subsection{Bangla-Capable Language Models}\label{bangla-capable-language-models}

\hl{BanglaBERT and BanglaT5 remain strong and efficient Bangla baselines
for classification and generation \cite{bhattacharjee2022banglabert}, \cite{bhattacharjee2023banglanlg}. TigerLLM continually
pre-trains small models on about 10 million tokens from 163 NCTB
textbooks (grades 6--12) and instruction-tunes them on 100,000 Bangla
instructions created from 500 human-written seed tasks; the paper
describes 1B (LLaMA 3.2) and 9B (Gemma 2) models and reports that they
outperform other open Bangla LLMs \cite{raihan2025tigerllm}. The released 1B checkpoint,
however, has exactly the architecture and tokenizer of Gemma-3-1B, which
is why this thesis pairs it with Gemma-3-1B. TituLLMs continually
pre-train Llama-3.2 models on about 37 billion Bangla tokens with an
extended tokenizer; the 3B model is a base (not instruction-tuned) model
\cite{titulm2025}. General multilingual open models are the alternative: Gemma 3
and Gemma 4 \cite{gemma4}, \cite{gemma3}, Qwen3 \cite{qwen3}, Llama 3.2 \cite{meta2024llama32} and
Tiny Aya, whose ``Fire'' variant targets South Asian languages \cite{tinyaya}.
Gemma 4 is trained on over 140 languages, is available in small
``effective 2B'' and ``effective 4B'' edge sizes with a context of up to
128K tokens, and is licensed under Apache 2.0 \cite{gemma4}. Whether Bangla
continued pre-training still helps once a model is fine-tuned for the
task had not been tested under a controlled budget; Chapter 5 does so.}

Evaluations of general LLMs on Bangla point the same way. BenLLM-Eval
tested ChatGPT, LLaMA-2 and Claude-2 zero-shot on a range of Bengali NLP
tasks. In a few tasks these models matched or exceeded fine-tuned
models, but in most tasks their performance was poor
\cite{kabir2024benllmeval}. Massively multilingual instruction tuning is
one response: the Aya model follows instructions in 101 languages, more
than half of them lower-resourced, and outperforms earlier multilingual
instruction-tuned models such as mT0 and BLOOMZ on most tasks
\cite{ustun2024aya}. On Bangla QA, LLMs still answer less accurately than
in English \cite{shafayat2024benqa}. Together, these results suggest
that an open multilingual model needs task-specific adaptation for
Bangla. They also motivate comparing several open base models under
one identical fine-tuning budget, rather than choosing one by
reputation.

\subsection{Bangla Healthcare and Education Applications}\label{bangla-healthcare-and-education-applications}

\hl{DocTalkBN is a corpus of expert telemedicine conversations from
national health programmes: 557.63 hours of audio, 1,515 multi-turn
patient conversations and 10,274 host--doctor question--answer
transcripts in 26 specialties, about 1.7 million tokens in total
\cite{saha2026doctalkbn}. Its benchmarks show that classifying doctor advice as safe or
harmful is easy for current models (macro-F1 above 0.92), while four-way
triage is hard (best macro-F1 0.463); the authors name
retrieval-augmented generation over the corpus as future work.
BanglaCHQ-Summ contains 2,350 real consumer health questions from an
online consultation platform with human-written summaries; fine-tuned
BanglaT5 reached a ROUGE-L of 48.35 \cite{khan2023banglachqsumm}. BanglaMedQA evaluates RAG
strategies for Bangla medical multiple-choice questions and finds large
gains from retrieval (for example 42.15\% \ensuremath{\rightarrow} 66.60\% for Llama-3.1-8B)
\cite{banglamedqa}, in line with English evidence that retrieval improves medical
QA \cite{xiong2024medrag} and that LLMs encode clinical knowledge \cite{singhal2023medpalm}.}

\hl{Medical entity recognition on such informal text is hard.
Bangla-HealthNER annotates 31,783 consumer health texts from a public
online health platform with seven entity types (symptom, health
condition, age, medicine, dosage, medical procedure and specialist), and
the best model, BanglishBERT, reaches an F1 of only 56.13\% \cite{khan2023nervous}.
The texts frequently switch between Bangla and English and come from
writers of different regions, and Bangla's rich inflection and flexible
word order make entity extraction harder than in English. Bangla MedER
annotates 6,895 medical statements with entities such as medicine and
chemical names, diseases, organs, pharmacological classes and hormones,
and reports 89.58\% accuracy with an ensemble of two BERT branches,
11.80 percentage points above a single BERT model \cite{aurpa2026banglameder}.}

\hl{For education, a cross-national survey of 318 university students
from the US and Bangladesh found that students already use LLMs for
learning, with LLM use correlated with perceived literacy benefits (r =
.59) and optimism towards AI (r = .41), alongside concerns about
over-reliance, ethics and critical thinking \cite{hossain2025ailiteracy}. This evidence,
however, concerns university students who are comfortable with English
and Standard Bangla, not school students who might benefit most from
curriculum support in their own dialect. Farzana et al. addressed the
grading of written answers, which is still done manually in remote and
rural schools with few qualified teachers: they built a synthetic
bilingual dataset of 122,710 questions with 490,391 question--answer
instances and fine-tuned Qwen3-8B with QLoRA to produce a score and
short feedback, which graded Bangla and English written answers close to
teachers (Spearman \ensuremath{\rho} = 0.936, MAE = 0.725 on 259 teacher-graded
instances) \cite{farzana2026semantic}. This shows that parameter-efficient fine-tuning of
open models is practical for Bangla education, and that answers should
be judged by meaning rather than word overlap. None of these systems
accepts dialect questions, retrieves evidence and generates grounded
answers.}

\subsection{Dialect Processing Beyond Bangla}\label{dialect-processing-beyond-bangla}

The dialect problem is not specific to Bangla. DIALECTBENCH collected
10 text-level tasks covering 281 varieties and found clear performance
gaps between standard and non-standard varieties of the same language
\cite{faisal2024dialectbench}. A study of LLMs on regional dialects in
machine translation and speech recognition found that the size of this
dialect gap depends on the training data, but not consistently across
models and languages, so that no single fix works everywhere
\cite{kantharuban2023dialectgap}. Arabic, which has a standard written
variety and many spoken dialects, is the best-studied case: the MADAR
corpus provides parallel sentences in 25 city dialects
\cite{bouamor2018madar}, and the NADI shared tasks have run dialect
identification at country and province level (NADI 2020 covered 100
provinces in 21 countries) \cite{abdulmageed2020nadi} and, more recently,
dialect-to-standard translation, where scores are still low
\cite{abdulmageed2023nadi}. Swiss German, which is spoken but rarely
written in a standard form, has been normalized into Standard German
before translation with about 60,000 parallel words, and character-based
translation worked best \cite{honnet2018swissgerman}.

Two lessons carry over to Bangla. First, the dialect gap must be
measured, because benchmarks in the standard variety hide it
\cite{faisal2024dialectbench}, \cite{kantharuban2023dialectgap}. Second,
normalizing dialect input into the standard variety lets existing
standard-variety systems be reused \cite{honnet2018swissgerman}. Most of
this work, however, evaluates identification or translation as end
tasks. This thesis measures the effect of normalization on a downstream
task, question answering, using the same questions in each variety.

\subsection{Combining and Routing LoRA Adapters}\label{combining-and-routing-lora-adapters}

Adapters are small trainable modules added to a frozen pre-trained model
\cite{houlsby2019adapters}, and LoRA is a low-rank form of them
\cite{hu2022lora}. Because each adapter is small, one base model can hold
many of them, which raises the question of how to combine them.
AdapterFusion learns to combine the representations of several
task adapters for a new target task and outperformed full fine-tuning
and multi-task learning on 16 language-understanding tasks
\cite{pfeiffer2021adapterfusion}. LoraHub composes several LoRA modules
for an unseen task by searching for mixing weights from a few examples,
and comes close to in-context learning without needing the examples in
the prompt \cite{huang2024lorahub}. Mixture of LoRA Experts (MoLE) learns a
gating function that weights the outputs of several LoRA modules in each
layer, and fuses them better than simple arithmetic merging
\cite{wu2024mole}. X-LoRA uses the model's hidden states to mix LoRA
experts at every layer and token \cite{buehler2024xlora}. On the systems side,
S-LoRA serves thousands of LoRA adapters from one base model, with up to
four times the throughput of earlier serving libraries
\cite{sheng2024slora}, which shows that many adapters on one shared model
is a practical way to deploy several skills.

These methods mix adapters \emph{inside} the network, with weights that
are learned and that change from token to token. This thesis takes a
simpler, system-level approach: a separate, inspectable CPU router
chooses \emph{one} whole adapter per step (the normalization expert, then
one domain expert), so every decision can be logged, evaluated and
corrected on its own. Because a learned mixture would be harder to
inspect and to test component by component, the thesis instead compares
routed experts with a single pooled adapter trained at equal compute.

\subsection{Re-ranking and Hybrid Retrieval}\label{re-ranking-and-hybrid-retrieval}

A retriever such as BM25 \cite{robertson2009bm25} or a dense encoder
\cite{karpukhin2020dpr}, \cite{chen2024bgem3} returns candidates
quickly, but its ranking is often imperfect. Re-ranking passes the top
candidates to a stronger model that reads the question and each passage
together. A BERT cross-encoder used in this way greatly improved passage
ranking \cite{nogueira2019passage}, and the sequence-to-sequence ranker
monoT5 was especially strong when training data were scarce and
transferred zero-shot to other collections \cite{nogueira2020monot5}.
ColBERT keeps much of this quality at far lower cost by comparing token
embeddings late in the computation, and runs two orders of magnitude
faster than BERT-based rankers \cite{khattab2020colbert}. Multilingual
cross-encoders such as bge-reranker-v2-m3 make re-ranking available for
Bangla \cite{bgereranker2024}. Rankings from several retrievers or
several query formulations can also be fused with Reciprocal Rank Fusion
\cite{cormack2009rrf}, as in RAG-Fusion \cite{rackauckas2024ragfusion}.
When the right collection is uncertain, a router can choose between
retrievers, as RouterRetriever does over a mixture of expert embedding
models \cite{lee2025routerretriever}. Bangla QA work has used retrieval
mainly with a single retriever \cite{banglamedqa}. This thesis chooses a
retriever per domain on validation data, and its answer-time extension
adds cross-encoder re-ranking and a search of a second knowledge base
when the domain router is unsure (Chapter~4).

\subsection{Grounding, Abstention and Dialect Robustness}\label{grounding-abstention-and-dialect-robustness}

\hl{Grounding answers in retrieved evidence is the standard remedy for
hallucination in knowledge-intensive tasks \cite{ji2023survey}, \cite{lewis2020rag}. Self-RAG
and Adaptive-RAG decide when and how to retrieve \cite{jeong2024adaptiverag}, \cite{asai2024selfrag},
SQuAD 2.0 and selective QA teach models to abstain \cite{rajpurkar2018squad2}, \cite{kamath2020selective},
and BanglaRQA and NCTB-QA provide real unanswerable Bangla questions for
learning this behaviour \cite{ekram2022banglarqa}, \cite{nctbqa}. For dialect robustness in
English, Multi-VALUE generates synthetic dialect variants for evaluation
and augmentation \cite{ziems2023multivalue}, and TADA trains a task-agnostic dialect
adapter that aligns dialect and standard representations; its gains are
larger for dialects far from the standard (+2.8 to +3.9 points) than for
dialects close to it (+0.3 to +0.4) \cite{held2023tada}. What is missing is a
system that treats uncertainty about the question itself, caused by
dialect and routing errors, as well as uncertainty about the evidence.}

{\small\setlength{\tabcolsep}{4pt}
\begin{longtable}{@{}H{0.135\textwidth}H{0.195\textwidth}H{0.165\textwidth}H{0.195\textwidth}H{0.170\textwidth}@{}}
\caption{Summary of the most closely related studies and the gap each leaves for this thesis.}\label{tab:related-studies}\\
\toprule
\textbf{Work} & \textbf{Task and data} & \textbf{Method} & \textbf{Key result} & \textbf{Gap} \\
\midrule
\endfirsthead
\toprule
\textbf{Work} & \textbf{Task and data} & \textbf{Method} & \textbf{Key result} & \textbf{Gap} \\
\midrule
\endhead
\bottomrule
\endfoot
Vashantor \cite{faria2023vashantor} & Dialect\ensuremath{\rightarrow}\allowbreak{}Standard translation and region detection; 5 regions, 12,500 pairs & BanglaT5, mT5; DialectBanglaT5/\allowbreak{}BERT (v2) & BLEU 44.03 (Ctg), 51.08 (Syl) with BanglaT5; 89.02\% region accuracy (v2) & Everyday sentences; no QA \\
Kothon, ONUBAD, ChatgaiyyaAlap \cite{kothon2026}, \cite{sultana2025onubad}, \cite{chatgaiyyaalap2025} & Dialect--Standard parallel corpora & Dataset papers & Human-validated pairs & No downstream task \\
5-Dialects-BN \cite{jawad2026fivedialects} & 6,000 utterances, 5 dialects, 5 aligned fields & 7 LLMs; prompting and LoRA & LoRA (160/\allowbreak{}dialect) 73.6 BLEU to English; some models below copy baseline & Normalization only prompted; no QA \\
Bridging Dialects \cite{khandaker2024bridging} & Standard\ensuremath{\rightarrow}\allowbreak{}dialect on Vashantor & BanglaT5, mT5, mBART-50 & WER 0.71 (Ctg) vs 0.24 (Mym) & Reverse direction; no QA \\
BanglaDialecto \cite{samin2024bangladialecto} & Noakhali speech standardization, 10 h & Whisper, BanglaT5, TTS & ASR CER 0.8\%; BLEU 41.6 & One dialect; no QA \\
Bangla\allowbreak{}CHQ-Prantik \cite{prantik2025} & Health questions: 2,350 Sylheti, 500 Chittagonian & 4 LLMs, prompting & Best BLEU 23.67 (Syl, 5-shot) & Standard\ensuremath{\rightarrow}\allowbreak{}dialect; no fine-tuning; no QA \\
Sylheti-CAP \cite{prama2025sylheticap} & Sylheti\ensuremath{\rightarrow}\allowbreak{}Standard, Vashantor test & Context-aware LLM prompting & Best chrF 53.72 & Prompting only; no QA \\
RAG dialect MT \cite{sami2025comparative} & Standard\ensuremath{\rightarrow}\allowbreak{}dialect, 6 dialects & RAG few-shot, 10 LLMs & BLEU 5--12 without examples; Chittagong 9\ensuremath{\rightarrow}\allowbreak{}26 with sentence-pair retrieval & Translation only \\
BanglaRQA \cite{ekram2022banglarqa} & Reading comprehension, 14,889 QA, 3,631 unanswerable & BanglaT5 and encoders & 62.42 EM, 78.11 F1 & Standard only; passage given \\
NCTB-QA \cite{nctbqa} & Textbook QA, 87,805 pairs, 42.75\% unanswerable & BERT, RoBERTa, ELECTRA & BERT F1 0.150\ensuremath{\rightarrow}\allowbreak{}0.620 & Standard only; passage given; model-generated questions \\
BEnQA, SOMAJGYAAN \cite{shafayat2024benqa}, \cite{shifa2025somajgyaan} & Bangla MCQ benchmarks & LLM prompting & Bangla below English; closed models lead & Standard only; multiple choice \\
bn\_squad \cite{mayeesha2021bengaliqa} & SQuAD 2.0 translated to Bangla; small human-annotated Wikipedia QA set & mBERT, DistilBERT, RoBERTa & Established machine-vs-human (children) baseline gap & Translated data is noisy; no retrieval-augmented answering \\
NCTB Multi-BERT QA \cite{khondoker2024unlocking} & NCTB passage QA, classes 6--10; about 3,000 pairs & BERT Base, Bangla-BERT, RoBERTa & Bangla-BERT F1 0.75, EM 0.53 & Passage given; Standard Bangla \\
BanglaCQA \cite{azmary2026banglacqa} & Counterfactual contexts for BanglaRQA & BanglaT5, LLMs & Models follow the passage & No retrieval or dialects \\
TigerLLM, TituLLMs \cite{raihan2025tigerllm}, \cite{titulm2025} & Bangla LLMs & Continued pre-training & Better than other open Bangla LLMs & No dialects; not compared after task fine-tuning \\
AmarDoctor \cite{nahar2025amardoctor} & Triage app, 185 vignettes & Knowledge graph, dialect symptom lexicon & 81.08\% top-1 diagnostic precision & Lexicon coverage; not QA \\
DocTalkBN \cite{saha2026doctalkbn} & Telemedicine corpus, 10,274 QA & LLM and encoder benchmarks & Triage macro-F1 0.463; safety F1 > 0.92 & No retrieval; no dialects \\
Bangla\allowbreak{}CHQ-Summ \cite{khan2023banglachqsumm} & 2,350 real consumer health questions with human summaries & Fine-tuned BanglaT5 for summarization & ROUGE-L 48.35\% & Summarization only, not QA; no answers provided \\
Bangla\-MedQA \cite{banglamedqa} & Bangla medical MCQ & 6 RAG strategies & 42.15\% \ensuremath{\rightarrow}\allowbreak{} 66.60\% with RAG & Standard only; MCQ \\
\bottomrule
\end{longtable}
}

\section{Summary of Findings and Research Gap}\label{summary-of-findings-and-research-gap}

\begin{itemize}
\item \hl{Regional varieties differ strongly from Standard Bangla, and
Chittagonian and Sylheti are among the hardest to model in both
translation directions \cite{jawad2026fivedialects}, \cite{faria2023vashantor}, \cite{khandaker2024bridging}, \cite{prantik2025}.}

\item \hl{Parallel dialect data exist (Vashantor, Kothon, ONUBAD,
ChatgaiyyaAlap, 5-Dialects-BN, BanglaCHQ-Prantik) and offer several
thousand dialect--Standard pairs for each target dialect, but they are
mostly everyday sentences, and medical dialect data are scarce,
especially for Chittagonian \cite{faria2023vashantor}, \cite{jawad2026fivedialects}, \cite{prantik2025}.}

\item \hl{Prompting is not enough for dialects: some LLMs fall below the
copy-input baseline on normalization, and prompted Standard\ensuremath{\rightarrow}dialect
translation of health questions reaches only modest BLEU \cite{jawad2026fivedialects},
\cite{prantik2025}. Small amounts of supervised adaptation give large gains
\cite{jawad2026fivedialects}, \cite{faria2023vashantor}.}

\item \hl{Retrieval helps both dialect handling and domain QA: retrieved
dialect--Standard pairs improve LLM dialect translation \cite{sami2025comparative}, and
retrieval improves medical QA \cite{banglamedqa}, \cite{xiong2024medrag}.}

\item \hl{Bangla QA benchmarks with unanswerable questions exist (BanglaRQA,
NCTB-QA), but all Bangla QA benchmarks, including BEnQA and SOMAJGYAAN,
are in Standard Bangla, and the earlier QA systems assume the correct
passage is given \cite{ekram2022banglarqa}, \cite{nctbqa}, \cite{shafayat2024benqa}, \cite{shifa2025somajgyaan}.}

\item \hl{LLMs are not as good as English on Bangla questions, and open
source models are behind closed source models \cite{shafayat2024benqa}, \cite{shifa2025somajgyaan}.}

\item \hl{Bangla-adapted LLMs exist \cite{raihan2025tigerllm}, \cite{titulm2025}, but their benefit
over general multilingual models after task fine-tuning had not been
measured under an equal budget.}

\item \hl{Bangla healthcare AI ranges from knowledge-graph triage to
telemedicine corpora \cite{nahar2025amardoctor}, \cite{saha2026doctalkbn}, \cite{khan2023banglachqsumm}; triage remains hard
for LLMs, and consumer health text is informal and code-mixed \cite{khan2023nervous},
so explicit safety rules and referral are needed.}

\item \hl{In Bangladesh, students are already employing
LLM\textquotesingle s in learning but the studied tools were in Standard
Bangla or English and were used by university level learners \cite{hossain2025ailiteracy}.}

\item \hl{When the passage contradicts pretraining knowledge, Bangla QA
models largely follow the passage \cite{azmary2026banglacqa}, which supports grounding
answers in retrieved evidence. The previous NCTB and Wikipedia QA
systems were extractive and assumed that the passage was given \cite{mayeesha2021bengaliqa},
\cite{khondoker2024unlocking}.}

\item \hl{Parameter-efficient fine-tuning makes it possible to adapt one
open base model to several tasks with small adapters on one GPU
\cite{hu2022lora}, \cite{dettmers2023qlora}. This enables a modular architecture in which the
same base model serves as the dialect expert and as the domain experts.}

\item Dialect gaps are found across languages, and normalizing into the
standard variety is an established strategy, but its effect on a
downstream task such as question answering is rarely measured
\cite{faisal2024dialectbench}, \cite{kantharuban2023dialectgap},
\cite{honnet2018swissgerman}.

\item Many LoRA adapters can share one base model and be served together
\cite{sheng2024slora}. Existing methods mix them inside the network with
learned weights \cite{wu2024mole}, \cite{huang2024lorahub}, while
explicit, inspectable routing between whole adapters has received less
attention.
\end{itemize}

\hl{Hence, this thesis topic, Dialect-Aware Bangla Question Answering
with Routed Domain Experts and Retrieval-Augmented Generation for
Healthcare and Education, targets three gaps which are yet to be
resolved in previous studies, which are outlined below:}

\begin{itemize}
\item \hl{\textbf{Dialect-aware QA.} Dialect identification and
normalization have been studied as end tasks, but not connected to
retrieval and answering: no system lets a Chittagonian or Sylheti
speaker ask a health or education question and receive a grounded
answer.}

\item \hl{\textbf{Grounded and safe answering in one pipeline.} Bangla QA
and LLM studies have not combined domain routing, retrieval, abstention,
clarification and medical safety rules in one evaluated system.}

\item \hl{\textbf{Per-variety evaluation.} Existing Bangla QA results are
reported only for Standard Bangla; the end-to-end gap between Standard
and dialect users, and how much of it each component closes, has not
been measured.}
\end{itemize}
